\let\negmedspace\undefined
\let\negthickspace\undefined
\documentclass[journal]{IEEEtran}
\usepackage[a5paper, margin=10mm, onecolumn]{geometry}
%\usepackage{lmodern} % Ensure lmodern is loaded for pdflatex
\usepackage{tfrupee} % Include tfrupee package

\setlength{\headheight}{1cm} % Set the height of the header box
\setlength{\headsep}{0mm}     % Set the distance between the header box and the top of the text
\usepackage{gvv-book}
\usepackage{gvv}
\usepackage{cite}
\usepackage{amsmath,amssymb,amsfonts,amsthm}
\usepackage{algorithmic}
\usepackage{graphicx}
\usepackage{textcomp}
\usepackage{xcolor}
\usepackage{txfonts}
\usepackage{listings}
\usepackage{enumitem}
\usepackage{mathtools}
\usepackage{gensymb}
\usepackage{comment}
\usepackage[breaklinks=true]{hyperref}
\usepackage{tkz-euclide} 
\usepackage{listings}
% \usepackage{gvv}                                        
\def\inputGnumericTable{}                                 
\usepackage[latin1]{inputenc}                                
\usepackage{color}                                            
\usepackage{array}                                            
\usepackage{longtable}                                       
\usepackage{calc}                                             
\usepackage{multirow}                                         
\usepackage{hhline}                                           
\usepackage{ifthen}                                           
\usepackage{lscape}



\usepackage{amsmath,amssymb}
\usepackage{booktabs}
\usepackage{tikz}
\usetikzlibrary{arrows.meta,angles,quotes}

\begin{document}

\bibliographystyle{IEEEtran}
\vspace{3cm}

\title{2.3.17}
\author{EE25BTECH11014 - Bhoomika Lokesh}
% \maketitle
% \newpage
% \bigskip
{\let\newpage\relax\maketitle}

\renewcommand{\thefigure}{\theenumi}
\renewcommand{\thetable}{\theenumi}
\setlength{\intextsep}{10pt} % Space between text and floats


\numberwithin{equation}{enumi}
\numberwithin{figure}{enumi}
\renewcommand{\thetable}{\theenumi}

\textbf{Question:} Find the properties of the triangle with vertices \brak{-4,0},\brak{0,3} and \brak{4,0}. 

\textbf{Solution:}
\begin{table}[H]    
  \centering
  \begin{tabular}[12pt]{ |c| c|}
    \hline
    \textbf{Name} & \textbf{Point}\\ 
    \hline
	$\vec{A}$ & \myvec{-4\\0}\\
    \hline 
	$\vec{B}$ & \myvec{0\\3}\\
    \hline
	$\vec{C}$ & \myvec{4\\0}\\
    \hline
    \end{tabular}
  \caption{Variables Used}
  \label{tab:2.3.17}
\end{table}

Vectors of sides of the triangle are,

\begin{align}
    \brak{\vec{A}-\vec{B}}=\myvec{-4\\0}-\myvec{0\\3}=\myvec{-4\\-3}
\end{align}

 \begin{align}
\brak{ \vec{B}-\vec{C}}=\myvec{0\\3}-\myvec{4\\0}=\myvec{-4\\3}
 \end{align}

 \begin{align}
 \brak{\vec{A}-\vec{C}}=\myvec{-4\\0}-\myvec{4\\0}=\myvec{-8\\0}
 \end{align}

To check for orthogonality, the dot product of the vectors of the sides of the triangle should be equal to zero.
\begin{align}
     \brak{\vec{A}-\vec{B}}^T \cdot \brak{ \vec{B}-\vec{C}}=\myvec{-4&-3}\cdot \myvec{-4\\3}=16-9=7\neq 0\\
     \brak{ \vec{B}-\vec{C}}^T \cdot  \brak{\vec{A}-\vec{C}}= \myvec{-4&3} \cdot \myvec{-8\\0}=32+0=32 \neq 0\\
      \brak{\vec{A}-\vec{C}}^T \cdot \brak{\vec{A}-\vec{B}}=\myvec{-8&0} \cdot \myvec{-4\\-3}=32-0=32 \neq 0
     \end{align}

\begin{center}
Hence, the given triangle is \textbf{NOT RIGHT ANGLED}.
\end{center}

To check if the given triangle is isosceles,a median should be perpendicular to a side of the triangle.\\
Let the midpoints of the sides \textbf{AB},\textbf{BC},\textbf be $\vec{x},\vec{y}$and $\vec{z}$.
\begin{align}
\vec{x}=\frac{\vec{A}+\vec{B}}{2}=\myvec{-4\\0}+\myvec{0\\3}=\myvec{-2\\\frac{3}{2}}\\ \vec{y}=\frac{\vec{B}+\vec{C}}{2}=\myvec{0\\3}+\myvec{4\\0}=\myvec{2\\\frac{3}{2}}\\
\vec{z}=\frac{\vec{A}+\vec{C}}{2}=\myvec{-4\\0}+\myvec{4\\0}=\myvec{0\\0} 
\end{align}
The vectors of medians of the triangle are:
\begin{align}
    \vec{x}-\vec{C}=\myvec{-2\\\frac{3}{2}}-\myvec{4\\0}=\myvec{-6\\\frac{3}{2}}\\
    \vec{y}-\vec{A}=\myvec{2\\\frac{3}{2}}-\myvec{-4\\0}=\myvec{6\\\frac{3}{2}}\\
    \vec{z}-\vec{B}=\myvec{0\\0}-\myvec{0\\3}=\myvec{0\\\frac{3}{2}}
\end{align}
\begin{align}
     \brak{\vec{x}-\vec{C}}^T \cdot \brak{\vec{A}-\vec{B}}=\myvec{-6&\frac{3}{2}} \cdot \myvec{-4\\-3}=24-\frac{9}{2}=\frac{39}{2} \neq 0\\
     \brak{ \vec{y}-\vec{A}}^T \cdot \brak{\vec{B}-\vec{C}}= \myvec{6&\frac{3}{2}} \cdot \myvec{-4\\3}=-24+\frac{9}{2}=\frac{-39}{2} \neq 0\\
     \brak{ \vec{z}-\vec{B}}^T \cdot \brak{\vec{A}-\vec{C}}= \myvec{0&\frac{3}{2}} \cdot \myvec{-8\\0}=0
\end{align}
 $\therefore$ We can conclude that triangle \textbf{ABC} is an \textbf{ISOSCELES} triangle.\\
 
 
See Fig. 0.1
 
\begin{figure}[H]
\begin{center}
\includegraphics[width=0.7\columnwidth]{Figs/fig_3.png}
\end{center}
\caption{Graph}
\label{Triangle ABC witn Medians}
\end{figure}


\end{document}